\fancyhead[C]{\textit{Solutions}}
\section*{\centering Solutions}

\subsection*{Solutions to T/F Questions}
	\begin{table}[!h]
		\centering
		\begin{tabular}{c@{\hspace{2.5em}}cp{0.75\textwidth}}
			true & false & \\
			\specialrule{\lightrulewidth}{0.25em}{1.25em}
			\dbubble & \bubble & Every matrix transformation \(T \colon \mathbb{R}^n \to \mathbb{R}^m\) defined by \(T(\vec{x}) = A\vec{x}\) 
						      is a linear transformation. \\[0.75em]
			\dbubble & \bubble & If \(\{\vec{v}_1, \vec{v}_2, \vec{v}_3\}\) is a linearly dependent set of vectors from \(\mathbb{R}^n\), the set
						      \(\{T(\vec{v}_1), \, T(\vec{v}_2), \, T(\vec{v}_3)\}\) is also linearly dependent for any linear transformation
						      \(T \colon \mathbb{R}^n \to \mathbb{R}^m\).\\[0.75em]
			\bubble & \dbubble & If \(T_A \colon \mathbb{R}^m \to \mathbb{R}^n\), with \(T(\vec{x}) = A\vec{x}\) then
						      \(A \in \mathbb{R}^{m \times n}\).\\[0.75em]
			\bubble & \dbubble & If \(T_A \colon \mathbb{R}^n \to \mathbb{R}^m\) with \(T(\vec{x}) = A\vec{x}\), then
						      the span of the columns of \(A\) is the codomain of \(T_A\).\\[0.75em]
			\dbubble & \bubble & If \(T\) is a linear transformation from \(\mathbb{R}^n\) to \(\mathbb{R}^m\) and there exists more than one vector 
						      \(\vec{v} \in \mathbb{R}^n\) such that \(T(\vec{v}) = \vec{b}\) for some  \(\vec{b} \in \mathbb{R}^m\), then there exists
						      a nontrivial solution to the equation \(T(\vec{x}) = \vec{0}\).\\[0.75em]
			\bottomrule
		\end{tabular}
	\end{table}

\subsection*{Short Answer Questions}
\begin{enumerate}
	\item \begin{enumerate}
			\item \(\mathbb{R}^3\).
			\item \(\mathbb{R}^2\).
			\item \(\vspan\left\{\begin{bmatrix}1 \\ 1\end{bmatrix}, \begin{bmatrix}2 \\ 9\end{bmatrix}, \begin{bmatrix}0 \\ 2\end{bmatrix}\right\}\). An
				alternative way to state this is to notice that the matrix has 2 pivots. Since it has a pivot every row, it's columns will span the whole of 
				\(\mathbb{R}^2\).
		\end{enumerate}
	\item Note that 
		\[\begin{bmatrix}1 \\ 0\end{bmatrix} = -1/2\begin{bmatrix}1 \\ 3\end{bmatrix} + 3/2\begin{bmatrix}1 \\ 1\end{bmatrix}, \qquad
		  \begin{bmatrix}0 \\ 1\end{bmatrix} = 1/2\begin{bmatrix}1 \\ 3\end{bmatrix} - 1/2\begin{bmatrix}1 \\ 1\end{bmatrix}.
		\] Then, 
		\[T\left(\begin{bmatrix}1 \\ 0\end{bmatrix}\right) 
		  = T\left(-1/2\begin{bmatrix}1 \\ 3\end{bmatrix} + 3/2\begin{bmatrix}1 \\ 1 \end{bmatrix}\right)
		  = -1/2 \cdot T\left(\begin{bmatrix}1 \\ 3\end{bmatrix}\right) + 3/2 \cdot T\left(\begin{bmatrix}1 \\ 1 \end{bmatrix}\right).\]
		Similarly, we get that 
		\[T\left(\begin{bmatrix}0 \\ 1\end{bmatrix}\right) 
		  = T\left(1/2\begin{bmatrix}1 \\ 3\end{bmatrix} - 1/2\begin{bmatrix}1 \\ 1 \end{bmatrix}\right)
		  = 1/2 \cdot T\left(\begin{bmatrix}1 \\ 3\end{bmatrix}\right) - 1/2 \cdot T\left(\begin{bmatrix}1 \\ 1 \end{bmatrix}\right).\]
		The standard matrix is therefore 
		\[\begin{bmatrix}T\left(\begin{bmatrix}1 \\ 0\end{bmatrix}\right) & T\left(\begin{bmatrix}0 \\ 1\end{bmatrix}\right) \end{bmatrix}
		  = \begin{bmatrix}1 & 0\\ -5/2 & 5/2\\ 1/2 & 1/2\end{bmatrix}\] The domain is \(\mathbb{R}^2\), and the range is the span of the columns of the
		standard matrix.
	\item This is actually \emph{impossible}. If, for any vector \(\vec{b} \in \mathbb{R}^m\), there exists a vector \(\vec{v} \in \mathbb{R}^n\) such that 
		\(T(\vec{v}) = \vec{b},\) then that means that matrix equation \(A\vec{x} = \vec{b}\) is consistent for any \(\vec{b} \in \mathbb{R}^m\). We know
		however that that's impossible since \(A\) in this case doesn't have a pivot in every row.

	\item It scales the \(x\)-component of a vector in \(\mathbb{R}^3\) and then projects it down onto the \(xy\)-plane.\footnote{In this special case, the operations are interchangeable. However, this is usually \emph{not} the case as you'll learn shortly.} The domain is \(\mathbb{R}^3\), and so is the codomain.
		The range is the span of the columns, or the set \[\left\{\!\left.\begin{bmatrix}b_1 \\ b_2 \\ b_3\end{bmatrix} \in \mathbb{R}^3 \;\right|\; b_3 = 0\!\right\}.\] They're equivalent.

	\item To scale the \(x\) component of a vector in \(\mathbb{R}^3\), the matrix \[\begin{bmatrix}2 & 0 & 0 \\ 0 & 1 & 0 \\ 0 & 0 & 1\end{bmatrix}\] is
		used. To project downwards onto the floor, we used the matrix \[\begin{bmatrix}1 & 0 & 0 \\ 0 & 1 & 0 \\ 0 & 0 & 0\end{bmatrix}\]. What happens
		if you multiplied these two matrices together? What does this tell you about matrix transformations composition?

	\item \[T\left(\begin{bmatrix}1 \\ 4 \\ 1\end{bmatrix}\right) = T\left(\begin{bmatrix}1 \\ 0 \\ 1\end{bmatrix}\right) 
		  \iff T\left(\begin{bmatrix}1 \\ 4 \\ 1\end{bmatrix}\right) - T\left(\begin{bmatrix}1 \\ 0 \\ 1\end{bmatrix}\right) = \vec{0}\]
		By the linearity of \(T\), this is the same as
		\[T\left(\begin{bmatrix}1 \\ 4 \\ 1\end{bmatrix} - \begin{bmatrix}1 \\ 0 \\ 1\end{bmatrix}\right) = \vec{0}.\] This shows that 
		\[\begin{bmatrix}1 \\ 4 \\ 1\end{bmatrix} - \begin{bmatrix}1 \\ 0 \\ 1\end{bmatrix}\] is a nontrivial solution to the equation \(T(\vec{x}) = \vec{0}.\)


\end{enumerate}